% Chapter 11 — The Settlement Interface. Author: KLEPPMANN. Budget 5 pp.
% Folds the settlement-obligation unit (its lifecycle node and the market-claim leg) into
% KLEPPMANN's voice. F3: SettlementState is not a free-standing type; the instructed-but-
% unsettled trade is a settlement-obligation unit whose instructed|settled|failed node lives
% in UnitStatus (Ch. 6). One product graph unifies market claim + fails cascade + settlement.
% Seam (for FORMALIS): MC-1 is created by the record-date contract firing, not by the
% settlement projection, which writes nothing. Numbers per the frozen thread map.

\section{The Settlement Interface}\label{ch:settlement}

This chapter discharges constitution clause C-Scope.9 (the settlement-projection layer, jointly with Chapter~\ref{ch:scope}); it builds on C-3.5 (every view a projection) and C-4.8 (what the record is).

\subsection{Settlement as a projection}

The settlement interface is where the closed ledger meets the open infrastructure
of settlement. It is a projection like any other (\S3): a function of the committed
log that reads the moves the ledger has admitted and emits \emph{settlement
instructions}. The division of labour is exact --- the ledger records \emph{what}
settles; the external infrastructure decides \emph{how}, and how is out of scope
(settlement finality, payment-system access, and central-securities-depository
connectivity are external authorities the ledger reconciles against, never
functions it performs). An instruction is the boundary object: it names the
transfer to be performed and carries none of the machinery that performs it.
Settlement in size may confirm in part: the ledger records the settled quantity each
confirmation carries, as a new external fact, while how a delivery is partialled,
tranched, or auto-split under settlement discipline belongs to the settlement layer ---
the \emph{how}, out of scope.

What the interface never does is write. It is not the Transaction Executor and
holds none of its privilege; it proposes nothing to the door and changes no
balance. The single-writer discipline (\S5) is not bent here --- the interface
reads and emits, and only the Transaction Executor writes. Every fact the interface
consumes is already on the record, and every fact it produces is an instruction,
never a move. Four properties follow, and each is load-bearing.

\begin{itemize}[nosep]
\item \emph{Pure.} It reads the recorded log alone --- no clock, no call to the
outside world --- so any party recomputes the same instructions from the record
(\S2, \S3).
\item \emph{Total.} It is defined on every committed transaction. A transaction
that carries no boundary-crossing move --- a moveless corporate-action component, a
holding in a virtual ledger (\S10) --- has the empty instruction list as its image,
by construction, not by a special rule.
\item \emph{Deterministic.} The same log yields the same instructions, always.
\item \emph{Idempotent.} Because it is a function of the recorded log and the
settlement-obligation unit's recorded lifecycle node, re-running it produces the same
instructions, and an instruction whose transfer the record already shows as settled is
recognised and not performed twice.
\end{itemize}

\begin{lstlisting}
settle :: Ledger -> [Instruction]   -- pure, total, deterministic, idempotent; writes nothing

data Instruction                    -- the boundary object: what settles, never how
  = Cash     WalletId WalletId UnitId MinorUnits  -- from, to, currency, amount
  | Delivery WalletId WalletId UnitId Quantity     -- from, to, security, quantity
\end{lstlisting}

One consequence is stated once and relied on throughout: \emph{settlement failure
is a recorded event, never a reversal}. A settlement that does not complete returns
as an external event, recorded as a new fact; the instruction's underlying move is
not undone. History is never rewritten (\S2): a trade ultimately unwound is unwound
by a compensating transaction, through the same door, never by deleting the booking
that instructed it.

\subsection{Four episodes}

\textbf{Dividend: cash projects to one instruction.} W-ALPHA holds $10{,}000$ ACME;
the declared dividend is $1.50$ per share, and nothing encumbers the holding. The
payment-date firing has already produced its transaction
(Chapter~\ref{ch:collateral}): the conditional payment obligation discharges
vacuously, and one move crosses the boundary --- owned(USD) \USD{15{,}000}, ACME's
issuer wallet $\to$ W-ALPHA.
\emph{The projection} reads that committed move.
\emph{The instruction:}
\begin{enumerate}[label=(\arabic*),nosep]
\item Cash \USD{15{,}000}, ACME issuer $\to$ W-ALPHA.
\end{enumerate}
The moveless component of the same transaction --- the obligation's creation and
vacuous discharge --- crosses nothing and projects to the empty list.
\emph{The door} is not the interface's: the move already passed the Transaction
Executor, and the projection only reads what was admitted.
\emph{The ledger determined and recorded the entitlement; the interface instructs
only the cash that leaves the ledger's edge.}

\medskip\noindent\textbf{Future: each day's margin, and the expiry unwind.} FUT-IDX
is settled-to-market, multiplier $10$; W-ALPHA is long one contract bought at
$995.00$. On each settlement print the future's contract fires and emits the day's
variation-margin move (Chapter~\ref{ch:collateral}); the interface projects each to
one instruction, and the expiry's final settlement to a settlement instruction that
also closes the contract.

\begin{center}
\begin{tabular}{@{}llll@{}}
\toprule
Recorded print & Variation margin & Committed move & Instruction \\
\midrule
day 1 \ $1{,}000.00$ & $+50$ & owned(USD) $50$, W-CCP $\to$ W-ALPHA & Cash $50$ to W-ALPHA \\
day 2 \ $990.00$ & $-100$ & owned(USD) $100$, W-ALPHA $\to$ W-CCP & Cash $100$ to W-CCP \\
day 3 \ $1{,}005.00$ & $+150$ & owned(USD) $150$, W-CCP $\to$ W-ALPHA & Settlement $150$ to W-ALPHA \\
\bottomrule
\end{tabular}
\end{center}

\noindent At expiry the position unwinds: the final settlement of \USD{150}
projects to the settlement instruction above, while the unit-state change that
retires the contract carries no move and projects to nothing.
\emph{The door} again belongs to the Transaction Executor; the interface reads
three admitted moves and emits three instructions.
Cumulative variation margin is $50-100+150 = 100$, the total profit and loss
(Chapter~\ref{ch:valuation}) --- but the interface never sees the profit and loss,
only the cash that crosses each day.
\emph{The ledger nets the day to a single move; the interface settles exactly that
move, and no valuation reaches the boundary.}

\medskip\noindent\textbf{Split: the moveless corporate action projects to nothing.}
ACME splits two-for-one, effective 2026-06-01. One transaction carries the
re-denomination and the new terms (Chapters~\ref{ch:objects} and~\ref{ch:marketdata}):
\begin{enumerate}[label=(\arabic*),nosep]
\item move owned(ACME) $10{,}000$, ACME issuer $\to$ W-ALPHA --- W-ALPHA's
$10{,}000$ shares become $20{,}000$;
\item a new ProductTerms version: price frame $100.00 \to 50.00$, with the market
data operators declared (spot scales by $\tfrac12$; the cash dividend figure scales
$1.50 \to 0.75$; proportional figures stand).
\end{enumerate}
\emph{The projection.} Component~(2) carries no move; its image is the empty list by
the projection's totality. Component~(1) is the ledger mirroring a re-denomination
the external register performs of its own accord as the corporate action; the
interface issues no delivery or payment instruction for it, because a split settles
nothing between counterparties.
\emph{The door} admitted a well-formed transaction whose consistency of reference
is checked in Chapter~\ref{ch:invariants} ($20{,}000$ against the fresh $50.00$,
never the stale $100.00$); the interface adds no check and emits nothing.
\emph{What crosses the boundary is a transfer; a corporate action that
re-denominates terms crosses nothing, so it projects to nothing --- by
construction.}

\medskip\noindent\textbf{Variance swap: the terminal settlement projects to one
instruction.} VS-1 is a one-year OTC variance swap on IDX, $252$ fixings, strike $400$
variance points at \USD{1{,}000} per point. Through its life each fixing is a
log-return the ledger computes from recorded IDX closes, not a transfer
(Chapters~\ref{ch:marketdata} and~\ref{ch:valuation}); none crosses the boundary,
and each projects to the empty list. At expiry the realised variance is $441$ points and the contract proposes one
settlement move --- owned(USD) $1{,}000\times(441-400) = \USD{41{,}000}$, from the
strike payer to the realised-variance receiver.
\emph{The projection} reads that one committed move.
\emph{The instruction:}
\begin{enumerate}[label=(\arabic*),nosep]
\item Cash \USD{41{,}000}, strike payer $\to$ realised-variance receiver.
\end{enumerate}
\emph{A whole year of fixings accrues one statistic and settles once; only the terminal
cash crosses the edge, and the interface instructs exactly it.}

\subsection{The settlement-obligation unit and the market claim}

One recorded fact upstream of the projection is load-bearing for its correctness:
the settlement level of a trade. A determination turns on it --- who receives a
distribution when a trade is instructed but unsettled across a record date --- and
the fact must be read, not assumed away. That fact is not floating free with nowhere
to live. Every right and obligation a legal contract creates is a unit (\S5.6), and an
instructed-but-unsettled trade is exactly an obligation: a delivery-versus-payment
obligation, the seller owing the security and the buyer the cash. So it is a unit ---
the \textbf{settlement-obligation unit}. Its lifecycle node --- \texttt{instructed},
then \texttt{settled} or \texttt{failed} --- is a property of the unit alone, the same
for every reader, so it lives in that unit's UnitStatus (Chapter~\ref{ch:homes}), which
is a home like any other. There is no free-standing settlement-state kept somewhere the
three homes do not reach: the routing below reads the unit's recorded node exactly as it
reads any other recorded fact, and sufficiency (Chapter~\ref{ch:homes}) holds with the
three homes, not a fourth.

\medskip\noindent\textbf{Owned re-books at instruction.} Title moves on the
instruction event, on the record (\S4), not on settlement confirmation --- the
trade-date booking of Chapter~\ref{ch:collateral}. The instant a trade is
instructed the owned plane shows the buyer; the settlement instruction the
interface projects is the promise that the external register will catch up. Between
the two lies the instruction-to-settlement gap, and the gap is carried explicitly
--- never elided, never back-filled at read time.

\medskip\noindent\textbf{A record date inside the gap.} Economic entitlement to a
distribution is fixed by the \emph{ex-date}, not by the record date alone. A trade
executed before the ex-date is \emph{cum} --- the price it paid still includes the
coming distribution --- and one executed on or after the ex-date is \emph{ex}, its
price already excluding it; the ex-date operator that separates the two frames is the
one Chapter~\ref{ch:marketdata} declares, and the cum/ex invariance witness that
forbids reading a price on the wrong side of it is catalogued in
Chapter~\ref{ch:invariants}. The routing predicate reads the trade's execution date
against that recorded ex-date; the record date only fixes the registration snapshot
the paying agent pays against. Whom the paying agent actually credits turns on a second
recorded fact --- whether the trade has \emph{settled} by the record date, so the
external register shows the buyer, or is still \emph{unsettled}, so it still shows the
deliverer. Two recorded facts, each binary --- cum or ex, settled or unsettled at the
record date --- make four quadrants. The paying agent pays the registered holder in
all four; in two the registered holder is the entitled party and nothing more is owed,
and in the other two the two diverge and a leg of the settlement-obligation unit
reconciles them.

\emph{The cum case.} A \emph{cum} sale still unsettled at the record date leaves the
owned plane showing the buyer --- who is cum-entitled --- while the paying agent
credits the \emph{record holder} at the external register, still the deliverer, the
seller. The two disagree by construction, and the disagreement is not a break to be
reconciled at report time: it is a recorded obligation. It is carried by the
settlement-obligation unit at its \texttt{instructed} node: across a record date in the
gap, the record-date firing walks that unit onto its market-claim branch, raising the
\textbf{market claim} --- the reconciling leg whose obligee holds it positive and whose
obligor holds it negative, the cum buyer's claim-for-equivalent on one side and the
deliverer's duty to remit on the other, discharged when the deliverer receives the
distribution.

\emph{The ex case.} An \emph{ex} sale unsettled at the record date raises no buyer
claim at all: the buyer bought ex, is not entitled to this distribution, and the
registered holder --- the deliverer --- keeps it, matching the ex-buyer's economic
expectation exactly. The owned plane shows the buyer here too, so owned read alone
would misroute; entitlement keys on the ex-date, and the record-date contract stays
silent.

\emph{The fourth quadrant: an ex sale settled before the record date.} The two cases
above both leave the trade unsettled at the record date. Settle it, and the register
catches up. For a \emph{cum} sale settled by the record date the register now shows
the buyer, who is cum-entitled --- registered holder and entitled party are the same,
and no claim arises. The last corner is the one the earlier enumeration missed. Take
an \emph{ex} sale --- the buyer bought without the coming distribution and is not
entitled --- that settles \emph{before} the record date, so the register shows the
buyer as the record holder. The paying agent credits the buyer, who is not entitled,
while the seller, who is, has been de-registered. This is the exact mirror of the cum
case: there the registered holder lagged the entitlement, here it leads it. The repair
is the mirror of the market claim --- the same settlement-obligation unit, its
market-claim leg raised at the \texttt{settled} node with the signs reversed: the
\textbf{mirror market-claim} from the wrongly-paid buyer to the entitled seller,
obligee the seller, obligor the buyer, discharged when the buyer receives the
distribution and remits it.

When does an ex sale settle before the record date at all? Under the ordinary calendar
it cannot: the ex-date is set relative to the record date and the settlement cycle
precisely so that a trade executed ex settles only \emph{after} the record date,
landing back in the ex-unsettled case. That coherence is an external market convention,
not a fact the ledger enforces; it is named TA-EXDATE in Chapter~\ref{ch:requirements}
and reconciled against at the boundary (Chapter~\ref{ch:marketdata}). It breaks
\emph{lawfully} under a due-bills regime for a large special dividend, whose ex-date
falls \emph{after} the payment date: there an ex sale can settle before the record
date, the fourth quadrant is reached, and the mirror market-claim leg routes the
distribution to the entitled seller --- on the record, by a leg the ledger raises, not
by a convention it assumes.

Here is the seam that keeps the projection pure: the market claim is created not by the
projection, which writes nothing, but by the contract that fires at the record date.
That contract reads the owned plane, the trade's cum/ex status against the recorded
ex-date, and the settlement-obligation unit's recorded lifecycle node, and proposes the
reconciling leg --- where the cum case calls for one --- as one component of the
determination transaction, which the Transaction Executor admits like any other. The
projection then reads the resulting committed moves; it originates none of them.

One product graph governs the settlement-obligation unit, and three mechanisms that were
once specified apart are its nodes and edges. The ordinary settlement is the spine,
\texttt{instructed}$\to$\texttt{settled}; the fails cascade (Chapter~\ref{ch:collateral},
\S9.3) is the edge \texttt{instructed}$\to$\texttt{failed}, on which the very same unit
\emph{becomes} the fails-cascade obligation; and the market claim is the branch raised at
the \texttt{instructed} node when a cum trade's record date falls in the gap, with its
mirror raised at the \texttt{settled} node when an ex trade's record date follows
settlement --- one mechanism, reflected across the settlement node. Being a product graph
like any other, its consistency is the four executable properties of
Invariant~\ref{inv:graph-consistency}, tested over generated settlement histories in
Chapter~\ref{ch:testability}; no new machinery is added, and no fourth home.

\begin{lstlisting}
-- The settlement-obligation unit (S5.6): an instructed-but-unsettled trade is a
-- delivery-versus-payment obligation, so it is a unit like any other.
data SettlementObligation = SettlementObligation
  { deliveryLeg   :: (WalletId, UnitId, Quantity)    -- the security the seller owes the buyer
  , paymentLeg    :: (WalletId, UnitId, MinorUnits)  -- the cash the buyer owes the seller
  , dueDate       :: Date
  , dischargePred :: Predicate  -- satisfied by settlement confirmation on the external register
  }
-- ONE product graph governs it (Ch. 3). Its lifecycle node lives in UnitStatus (Ch. 6):
--   instructed --confirmed--> settled     the ordinary walk; the unit retires unless a
--                                          record date straddles the settled node (see below)
--   instructed --fails------> failed      the fails-cascade obligation (S9.3): the SAME unit, later node
--   at INSTRUCTED, a CUM trade across a record date raises the market-claim leg (deliverer -,
--     cum-buyer +), valued at par, discharged by the deliverer's receipt of the distribution.
--   at SETTLED, an EX trade across a record date raises the MIRROR market-claim leg, signs
--     reversed (ex-buyer -, seller +): the fourth quadrant, discharged by the buyer's receipt
--     and remittance to the seller. Same mechanism, reflected across the settlement node.
-- Routing reads this node from UnitStatus; there is no free-standing settlement-state type.
\end{lstlisting}

\begin{proposition}[Entitlement routing]
The recipient of a distribution is a total function of the recorded owned plane, the
trade's cum/ex status against the recorded ex-date, and the settlement-obligation unit's
recorded lifecycle node at the record date. The last two facts are each binary --- cum or
ex, settled or unsettled at the record date --- and partition into four quadrants; the
paying agent credits the registered holder in all four, and a market-claim leg reconciles
exactly the two in which the registered holder is not the entitled party.
\emph{(i) Cum, unsettled:} the trade is still at its \texttt{instructed} node, so the
register shows the deliverer, and the settlement-obligation unit's market-claim leg
carries the reconciling claim from the deliverer to the cum-buyer.
\emph{(ii) Cum, settled:} the register shows the buyer, who is entitled, so no leg arises.
\emph{(iii) Ex, unsettled:} the register shows the deliverer, who is entitled, so no leg
arises.
\emph{(iv) Ex, settled --- the fourth quadrant:} the register shows the buyer, who is
\emph{not} entitled, so the same settlement-obligation unit, at its \texttt{settled} node,
raises the \emph{mirror} market-claim leg, signs reversed, carrying the reconciling claim
from the wrongly-paid buyer to the entitled seller.
No distribution reaches any wallet except by a recorded move against a determined
entitlement or the discharge of a market-claim leg --- direct or mirror --- of the
settlement-obligation unit.
\end{proposition}

\medskip\noindent\textbf{Worked micro-example (dividend).} Stock ACME pays a cash
dividend of $1.50$ per share; the record date is 2026-05-15. Wallet W-ALPHA holds
$10{,}000$ ACME and sells them to a buyer, wallet W-BETA: the sale is instructed on
2026-05-14 for settlement on 2026-05-16, so the record date falls in the gap. The
2026-05-14 execution precedes the ex-date associated with the 2026-05-15 record date,
so the sale is \emph{cum} --- which is why the distribution routes to the buyer; an
otherwise identical trade executed on or after the ex-date --- and, as here, still
unsettled at the record date --- would be \emph{ex} and raise no market claim. (Had that
ex trade instead settled before the record date, the fourth quadrant, it would raise the
mirror market-claim to the seller. The cash consideration settles symmetrically and does
not bear on the entitlement.)

\emph{2026-05-14 --- the sale; owned re-books at instruction.} One transaction:
\begin{enumerate}[label=(\arabic*),nosep]
\item move owned(ACME) $10{,}000$, W-ALPHA $\to$ W-BETA --- re-booked at
instruction; the settlement instruction, due 2026-05-16, is recorded as unsettled.
\end{enumerate}
The owned plane now shows W-BETA holding $10{,}000$ ACME and W-ALPHA holding none.

\emph{2026-05-15, record date --- determination and the market claim.} The
record-date watch fires; the stock's contract determines the entitlement and, reading
the trade's cum status against the ex-date and the settlement-obligation unit's recorded
node, proposes the reconciling leg. One transaction:
\begin{enumerate}[label=(\arabic*),nosep]
\item determination reads the owned plane --- W-BETA owns $10{,}000$ --- and records
the entitlement, $10{,}000 \times 1.50 = \USD{15{,}000}$, to W-BETA in PositionState;
\item the sale is cum and its settlement-obligation unit is still at \texttt{instructed}
across the record date, so the market-claim branch is walked: the paying agent will credit
the record holder, the deliverer W-ALPHA. Raise the market-claim leg MC-1 --- the
settlement-obligation unit on its record-date branch ---
valued at par: move owned(MC-1) $1{,}500{,}000$ (minor units, $=\USD{15{,}000}$),
W-ALPHA $\to$ W-BETA --- W-ALPHA holds $-1{,}500{,}000$ (the duty to remit once
received), W-BETA holds $+1{,}500{,}000$ (the claim-for-equivalent). MC-1's discharge
predicate is the deliverer's receipt; the obligation is contingent on it and is
matched by the deliverer's record-holder receivable, so the conduit's owned value is
neutral. While W-ALPHA performs the claim is valued at par, and it decouples to a
recovery claim on W-ALPHA's default.
\end{enumerate}
This transaction is admitted through the door; MC-1 exists because the record-date
firing wrote it, not because the projection did.

The sale settles on 2026-05-16, recorded as an external confirmation; this updates
the register but not the record-date snapshot that governs the distribution --- the
paying agent still credits the record holder as of 2026-05-15.

\emph{2026-05-22, payment date --- the distribution and the discharge.} The
payment-date watch fires. One transaction, two moves:
\begin{enumerate}[label=(\arabic*),nosep]
\item move owned(USD) \USD{15{,}000}, ACME issuer wallet $\to$ W-ALPHA --- the paying
agent credits the record holder;
\item move owned(USD) \USD{15{,}000}, W-ALPHA $\to$ W-BETA --- the market claim
discharges under its predicate, and MC-1 retires at the zero vector.
\end{enumerate}
W-ALPHA nets to zero across the two moves --- a conduit, not an owner; W-BETA ends
with the \USD{15{,}000} it economically owns. \emph{The door}: conservation per
(unit, coordinate) on every paired-leg move; MC-1 retires only from the zero vector;
idempotence --- a duplicate payment-date event finds MC-1 discharged and proposes
nothing. Each of these two cash moves is what the interface then projects to a cash
instruction; the market claim itself never crosses the boundary --- it is the
internal record that keeps the two boundary-crossing payments consistent.

\medskip\noindent\textbf{Settlement failure is a recorded event, never a reversal.}
Should the sale fail to settle on its due date, the settlement-obligation unit walks
to its \texttt{failed} node --- and the instruction's owned re-booking is not undone
and the market claim is not withdrawn. The walk has exactly one writer, and it is
named: the settlement-obligation unit's own smart contract, the one legal writer of
that unit's lifecycle node (Chapter~\ref{ch:homes}, Invariant~\ref{inv:one-writer}).
It proposes the transition on one of two recorded triggers --- an external fail
notice, admitted through the one door as a moveless observation-recording transaction
(Chapter~\ref{ch:marketdata}), or the due-date watch firing and finding the cumulative
confirmed settled quantity short of the instructed quantity by the deadline. The
second trigger reads quantity, not mere presence, so a partial settlement --- some
delivered, a residual failed --- fires it on the unmet residual rather than passing
silently; the quantity-aware splitting of the settlement-obligation unit into a
settled leg and a failed residual leg is a named open problem (Chapter~\ref{ch:scope}),
and until it is settled the residual is carried on the single \texttt{failed} node with
its recorded shortfall quantity. In both cases the Transaction Executor admits the
transition like any other transaction; no other machine writes the node, and no
\texttt{failed} node exists without one of the two triggers on the record. The settlement-obligation unit's market-claim leg carries
a deadline, a discharge predicate, and a fallback (\S14.1); by its deadline it is
discharged, compensated, or declared defaulted --- a failed distribution is a
visible, overdue item, never a silent balance edit.

\medskip\noindent\textbf{Load-bearing, not a reporting axis.} The settlement-obligation
unit's lifecycle node is therefore a correctness input, not a presentation dimension. A
design that read owned alone would instruct the venue to pay the buyer directly --- an
instruction it cannot honour, since the buyer is not yet the record holder --- and the
cash would arrive at the deliverer with nowhere on the record to go, resurfacing as
exactly the reconciliation break the ledger exists to abolish. Reading the
settlement-obligation unit's recorded node, and carrying the market claim on that unit,
makes the routing both faithful to what the venue will do and reproducible by any party
from the record alone (\S2, \S3). The construction is symmetric across the settlement
node, and that symmetry is the fourth quadrant: where a cum trade unsettled at the record
date raises the market claim from deliverer to buyer, an ex trade settled by the record
date raises the \textbf{mirror market-claim} on the same settlement-obligation unit, its
signs reversed, from the wrongly-paid buyer to the seller --- one mechanism, total over
all four quadrants of the gap.

\emph{The ledger says what settles and, across the gap, who is owed; the interface
projects that settled activity into instructions, writing nothing and inventing
nothing.}
